% Plantilla básica para ejemplar formal oficial de tesis de la Escuela Internacional de Doctorado (EIDUM) 
% de la Universidad de Murcia, España. (ejemplar formal)
% Antonio Maurandi López, 8 de marzo de 2023
%
% Notas: 
% - el dir "estilos" contiene un fichero variables.tex donde poner titulo, autor, etc... y otros de librerías etc...
% - en el dir "capitulos" van los capítulos a inc con un input
% - otro dir de "capitulos/tablas" para editarlas tranquilamente e inc con un input tb.
% - en el directorio "pics" pongo las imágenes
% - la biblio no es APA, pq no me gusta, tampoco es obligatorio q así lo sea.
% -
% ---------------------------------------------------------

\documentclass[a4paper,12pt]{book}
\input{./estilos/estilo-basico.tex} 
\input{./estilos/estilo-capitulos.tex}
\input{estilos/variables}

\begin{document}
%-------------------------------------------------------------------
%	PORTADA
%-------------------------------------------------------------------
\pagestyle{empty}
\input{capitulos/portada.tex}
\newpage 
\ % pagina ne blanco
\newpage
% CONTRAPORTADA con nombres de directoras
\input{capitulos/portada-interior.tex}
%\newpage

%-------------------------------------------------------------------
%	Frase Culterana
%-------------------------------------------------------------------
%\pagestyle{empty}
\frontmatter
\chapter*{}
\input{./capitulos/fraseculterana.tex}
%-------------------------------------------------------------------
%	DEDICATORIA
%-------------------------------------------------------------------
%\pagestyle{empty}
\frontmatter
\chapter*{}
\input{./capitulos/dedicatoria.tex}
%-------------------------------------------------------------------
%	AGRADECIMIENTOS
%-------------------------------------------------------------------
\chapter*{Agradecimientos}
%\markboth{AGRADECIMIENTOS23}{AGRADECIMIENTOS} % encabezado 
\input{./capitulos/agradecimientos.tex}

%-------------------------------------------------------------------
%	PREFACIO
%-------------------------------------------------------------------
\chapter*{Prefacio}
\input{./capitulos/prefacio.tex}

%-------------------------------------------------------------------
%	RESUMEN y SUMMARY para la versión oficial
%-------------------------------------------------------------------
\chapter*{Resumen y Summary}
\input{./capitulos/resumenYsummary.tex}

%-------------------------------------------------------------------
%	TABLA DE CONTENIDOS
%-------------------------------------------------------------------
%\setcounter{tocdepth}{1}
%level -1: part, 0: chapter, 1: section, etc.
\renewcommand{\baselinestretch}{0.1}\normalsize
\tableofcontents
\renewcommand{\baselinestretch}{0.3}\normalsize
\listoffigures
\renewcommand{\baselinestretch}{0.3}\normalsize
\listoftables
\renewcommand{\baselinestretch}{1.2}\normalsize

%-------------------------------------------------------------------
%	TESIS
%-------------------------------------------------------------------
\mainmatter %empieza la numeración de las páginas
\pagestyle{headings}

\def\tablename{Tabla}% para que denomine tabla a las tablas y no "cuadro"
%\doublespacing
\onehalfspace
%\singlespace
%\spacing{1.5}	
\setlength{\parskip}{3mm}% separación entre párrafo
\setlength{\parindent}{1cm}%identación primera linea de párrafo


\chapter{Introducción y antecedentes}
\input{./capitulos/10-intro.tex}

\chapter{Justificación y objetivos}
\input{./capitulos/20-just-objetivos.tex}

\chapter{Material y métodos}\label{Materialymetodos}
\input{./capitulos/30-matYmet.tex}

\chapter{Resultados y discusión}
\input{./capitulos/50-Resultados-y-Discusion.tex}

\chapter{Conclusiones}\label{Conclusiones}
\input{./capitulos/40-Conclusiones.tex}


%-------------------------------------------------------------------
%	BIBLIOGRAFÍA
%-------------------------------------------------------------------
%\backmatter
%\nocite{*}   %Ojo! si pones nocite{*} te incluye en bilbiografia todas las entradas del .bib, incluso cunado no la citas intext!!
%\bibliographystyle{unsrt} % igual q plain pero numera en orden de aparición
\bibliographystyle{apalike}

%\usepackage{apacite} %esto entes de begin document
%\DefineBibliographyStrings{spanish}{ andothers = {et\addabbrvspace al\adddot}, and = {y}, }
%\bibliographystyle{apacite}

%\bibliographystyle{plainnat}
\bibliography{referenciasBibliograficas.bib}


%-------------------------------------------------------------------
%	APÉNDICES
%-------------------------------------------------------------------

\addtocontents{toc}{\vspace{2em}} % Agrega espacios en la toc
\appendix % Los siguientes capítulos son apéndices

\chapter{Esto será un apéndice}\label{Apendice_A}
\input{./capitulos/110-apendice.tex}

\chapter{Esto será otro apéndice}\label{Apendice_B}
\input{./capitulos/120-apendice.tex}

\chapter{Esto será otro apéndice que es un PDF}\label{Apendice_C}
\includepdf[pages=-,templatesize={120mm}{175mm},noautoscale=false]{./capitulos/otrosDocs/documento-pdf-de-ejemplo.pdf}

\chapter{Esto será otro apéndice que es un PDF}\label{Apendice_C}
\includepdf[pages=-,templatesize={120mm}{175mm},noautoscale=false]{./capitulos/otrosDocs/codigo.pdf}

\newpage %añado una página adicional blanca
\thispagestyle{empty} % empty
\mbox{}

\addtocontents{toc}{\vspace{2em}} % Agrega espacio en la toc
\end{document}









